\section{Objetivos}

\subsection{Objetivo general}

Realizar la calibración de instrumentos de medición de longitud empleando patrones
certificados y procedimientos normalizados, y cuantificar la incertidumbre de las
mediciones mediante la ley de propagación de errores.

\subsection{Objetivos específicos}

\begin{enumerate}
  \item Aplicar los procedimientos PC-012, PC-014 y PC-029 del INACAL en la
        calibración del pie de rey, del comparador y del medidor de espesores por
        ultrasonido~\parencite{inacal_pc012, inacal_pc014, inacal_pc029}.
  \item Evaluar la conformidad del pie de rey y del comparador con su EMPE
        declarado, aplicando la regla de decisión
        ILAC-G8~\parencite{ilac2019decision}.
  \item Cuantificar la incertidumbre del pie de rey e identificar las componentes
        que la dominan.
  \item Cuantificar la incertidumbre del comparador e identificar las componentes
        que la dominan.
  \item Cuantificar la incertidumbre del medidor de ultrasonido y evaluar la
        influencia del patrón escalonado.
  \item Describir los fundamentos y equipos de los ensayos no destructivos
        (PT, MT, UT y VT) utilizados en la empresa.
\end{enumerate}
